\begin{trapbox}
	\begin{description}[style=nextline, leftmargin=2.8em, labelsep=0.5em, itemsep=0.35em]
		\item[\textbf{\textcolor{CTrap}{易错点一（三点共线）}}]若 $A,B,C$ 共线，则无法构成平行四边形，使用公式仍会计算出 $D$，但该 $D$ 无效。
		\item[\textbf{\textcolor{CTrap}{正确理解（先判共线）：}}]必须先验证 $A,B,C$ 不共线，再套用公式。
		\item[\textbf{\textcolor{CTrap}{错因分析（忽略条件）：}}]误以为任意三点都能确定平行四边形，未理解平行四边形定义。
	\end{description}
	\medskip
	\begin{description}[style=nextline, leftmargin=2.8em, labelsep=0.5em, itemsep=0.35em]
		\item[\textbf{\textcolor{CTrap}{易错点二（对角配对错误）}}]在套用公式时，错误地将 $A,B$ 当作对角顶点，而实际要求 $A,C$ 为对角，导致 $D$ 坐标错误。
		\item[\textbf{\textcolor{CTrap}{正确理解（明确对角）：}}]需根据题意或假设确定哪两点是对角顶点，再代入对应公式。
		\item[\textbf{\textcolor{CTrap}{错因分析（混淆顶点）：}}]混淆了平行四边形中对角线的概念，不知道哪些顶点是对角关系。
	\end{description}
	\medskip
	\begin{description}[style=nextline, leftmargin=2.8em, labelsep=0.5em, itemsep=0.35em]
		\item[\textbf{\textcolor{CTrap}{易错点三（漏解情况）}}]求平行四边形顶点时，只写出一个 $D$ 坐标，忽略其他两种可能。
		\item[\textbf{\textcolor{CTrap}{正确理解（三种可能）：}}]已知三个不共线点，第四个顶点有三种可能，需全部考虑（除非题目指定）。
		\item[\textbf{\textcolor{CTrap}{错因分析（思维不全面）：}}]没有考虑顶点顺序或对角线的多种选择。
	\end{description}
\end{trapbox}
